Generated GPU kernels require evidence of correctness and stable speed, not plausible source alone. We present \okf{}, an artifact-preserving evaluation harness for LLM-generated Triton kernels that combines static policy checks, correctness verification, CUDA-event timing, cache-state perturbation, repeated measurement sessions, session-level uncertainty summaries, and eager and compiler baselines. On a controlled internal fused8 suite, Gemini and OpenAI mini outperform \compile{} at median, by \speed{1.863} and \speed{1.835}, while remaining below PyTorch eager at median; deterministic templates provide the strongest overall floor. A deterministic \code{bias\_relu} candidate also changes from a \speed{1.029} single-run win to a \speed{0.976} repeat median, showing that measurement instability is not specific to model sampling. We additionally preserve a capped KernelBench L1 pilot and one repair pass. A post-hoc adapter audit found that those historical runs replaced the official \code{ModelNew} lifecycle with a free-function contract and reconstructed reference modules inside timed calls. We therefore retain their artifacts as a failure-characterization case study but exclude their correctness rates and speedups from supported performance claims pending corrected revalidation. The contribution is methodological: generated-kernel evaluation should distinguish correctness, compiler speedup, eager speedup, and repeat-stable speedup, while preserving enough evidence to audit the evaluator itself. We make neither a state of the art claim nor a benchmark-wide KernelBench claim.
